% --- MODELLAZIONE E DATASET ---
\section{Modeling and data generation}
%\subsection{Double Pendulum Setup}
\begin{frame}{Double Pendulum Setup}

    \begin{columns}
        \begin{column}{0.4\textwidth}
            \centering
            \begin{figure}
                \includegraphics[width=0.7\linewidth]{double_pendulum.png} 
                \caption{Double Pendulum Model.}
                \label{fig:double_pendulum}
            \end{figure}
        \end{column}
        
        \begin{column}{0.4\textwidth}
            \begin{block}{State Variables}
                
                \begin{equation}
                    \mathbf{q} = 
                    \begin{bmatrix}
                            q_1 \\
                            q_2 \\
                            \dot{q}_1 \\
                            \dot{q}_2 \\
                    \end{bmatrix}
                    = 
                    \begin{bmatrix}
                        \theta_1 \\
                        \theta_2 \\
                        \dot{\theta}_1 \\
                        \dot{\theta}_2 \\
                    \end{bmatrix}
                    \in \mathbb{R}^4
                \end{equation}
                
                Where:
                \begin{itemize}
                    \item  \[
                    \begin{bmatrix}
                            q_1 \\
                            q_2 \\ 
                    \end{bmatrix}
                    \in \mathbb{R}^2
                    \]   
                    are the angular positions.
                    \item \[
                        \begin{bmatrix}
                                \dot{q}_1 \\
                                \dot{q}_2 \\
                        \end{bmatrix}
                        \in \mathbb{R}^2
                    \]
                    are the angular velocities.

                \end{itemize}
            \end{block}
        \end{column}
    \end{columns}
\end{frame}

\begin{frame}{Double Pendulum Dynamics}

    From the Euler-Lagrange equations:
    \begin{equation}
        \frac{d}{dt} \left( \frac{\partial \mathcal{L}}{\partial \dot{q}_i} \right) - \frac{\partial \mathcal{L}}{\partial q_i} + \frac{\partial \mathcal{R}}{\partial \dot{q}_i}= \tau_i, \quad i = 1, 2
    \end{equation}
    where $\mathcal{L} = T - V$ is the Lagrangian. 
    \begin{itemize}
        \item $T$ is the kinetic energy.
        \item $V$ is the potential energy.
        \item $R$ is the Rayleigh dissipation function.
        \item $\tau_i$ are the generalized forces (torques) applied at the joints.
    \end{itemize}

    \vspace{10px}
    The dynamic model of the double pendulum can be expressed in the following matrix form.\footnote{Explicit expression reported in the appendix \ref{subsec:appendixA}.}

    \begin{block}{Dynamics in Matrix Form}
        %\text{It can be expressed as:}
        \begin{equation}
            M(\mathbf{q})\ddot{\mathbf{q}} + \mathbf{c}(\mathbf{q}, \dot{\mathbf{q}}) + D\dot{\mathbf{q}} + G(\mathbf{q}) = \mathbf{\tau}
        \end{equation}

    \end{block}

\end{frame}

%\subsection{Physic Model} 

\begin{frame}{Data Generation}

    
    Integration:
    \begin{itemize}
        \item \textbf{Forward Euler:}
        \begin{itemize}
            \item[\textcolor{ForestGreen}{\checkmark}] simple
            \item[\textcolor{red}{\times}] less accurate
        \end{itemize}
        \item \textbf{Runge-Kutta 4th Order (RK4):}
            \begin{itemize}
                \item[\textcolor{ForestGreen}{\checkmark}] more accurate
                \item[\textcolor{red}{\times}] computationally expensive
            \end{itemize}

    \end{itemize}

    \vspace{1px}

    \begin{block}{Simulation Setting}
        \begin{itemize}
        \item Duration $T = 10\,\text{s}$
        \item Step size $\Delta t = 0.01\,\text{s}$
        \item Initial state $q_0 \sim \mathcal{U}[-\pi, \pi], \dot{q_0} \sim \mathcal{U}[-2,2]$
        \item Physical parameters $m_1, m_2, l_1, l_2, g \text{(gravity)}, b_1, b_2 \text{(friction)}$
        \end{itemize}
    \end{block}
    

    \vspace{1px}

    \begin{block}{Generated datasets}
        \begin{itemize}
            \item[\color{blue}\blacktriangleright] \textcolor{blue}{\textbf{No friction}} $b_1 = b_2 = 0$
        
            \item[\color{orange}\blacktriangleright] \textcolor{orange}{\textbf{Low friction}} $b_1 = b_2 = 0.3$
            
            \item[\color{red}\blacktriangleright] \textcolor{red}{\textbf{High friction}} $b_1 = b_2 \approx 3.0$ \quad \small \textit{\textcolor{red}{Excluded for rapid energy dissipation}}
        \end{itemize}
    \end{block}

\end{frame}

% --- Slide 2: Dynamics ---


\begin{frame}{Motivation for Learning Dynamics}
    Why do we need data-driven approaches?
    
    \vspace{1em}

    \begin{itemize}
        \setlength\itemsep{1.5em} % Aumenta lo spazio tra i punti
        \item \textbf{Complexity}: real physical systems are chaotic and highly non-linear
        \item \textbf{Cost}: accurate parameter estimation via sensors is expensive and prone to noise
        \item \textbf{Adaptability}: analytical models fail when system properties change over time (e.g., wear and tear)
    \end{itemize}
    
    \vfill % Spinge il contenuto verso l'alto bilanciando lo spazio vuoto in basso
\end{frame}

  

